\chapter{点集拓扑}
\label{chap:md-03}

%——————————————————————————————————%

本章研究“空间”最基础的结构。集合只告诉我们有哪些点，拓扑则进一步规定哪些点彼此邻近、哪些集合可以称为开集、哪些映射保持连续变化。

本章的主线是
\[
\begin{aligned}
&\text{拓扑}
\longrightarrow
\text{子空间与积空间}
\longrightarrow
\text{连续映射与同胚}\\
&\longrightarrow
T_1,\ \text{Hausdorff}
\longrightarrow
\text{紧致、仿紧}
\longrightarrow
\text{连通、道路连通}
\longrightarrow
\text{同伦与基本群}.
\end{aligned}
\]

这些概念分别回答：

\begin{itemize}
    \item 什么叫一个点的邻域？
    \item 什么叫一个映射保持连续运动？
    \item 不同点能否被拓扑结构区分？
    \item 无穷多个局部区域能否由有限个控制？
    \item 空间能否被分成互不连通的部分？
    \item 一条闭合路径能否连续缩成一个点？
\end{itemize}

第四章将在此基础上加入局部欧氏性、Hausdorff性和第二可数性，定义拓扑流形。

%——————————————————————————————————%

\section{拓扑空间}

\subsection{开集与拓扑}

\begin{definition}[拓扑]
设$M$是一个集合，$\mathcal O\subseteq\P(M)$。若满足：
\begin{enumerate}
    \item
    \[
        \varnothing\in\mathcal O,\qquad M\in\mathcal O;
    \]
    \item 对任意指标集$A$和任意族$\{U_\alpha\}_{\alpha\in A}\subseteq\mathcal O$，
    \[
        \bigcup_{\alpha\in A}U_\alpha\in\mathcal O;
    \]
    \item 对任意有限个$U_1,\ldots,U_k\in\mathcal O$，
    \[
        \bigcap_{i=1}^kU_i\in\mathcal O;
    \]
\end{enumerate}
则称$\mathcal O$为$M$上的一个拓扑。$\mathcal O$中的集合称为开集。
\end{definition}

最容易被忽略的是并集和交集条件的不对称性：
\[
    \boxed{\text{任意并仍开，但只有有限交保证仍开。}}
\]

无限交一般不保持开性。例如在$\R$的标准拓扑中，令
\[
    U_n=\left(-\frac1n,\frac1n\right).
\]
每个$U_n$都是开集，但
\[
    \bigcap_{n=1}^{\infty}U_n=\{0\}
\]
不是$\R$中的开集。

\begin{definition}[拓扑空间]
集合$M$配备拓扑$\mathcal O$后，称$(M,\mathcal O)$为一个拓扑空间，通常简写为$M$。
\end{definition}

同一个集合可以配备不同拓扑，从而得到不同的空间结构。最粗的拓扑和最细的拓扑分别是：
\[
    \mathcal O_{\mathrm{ind}}=\{\varnothing,M\},
    \qquad
    \mathcal O_{\mathrm{disc}}=\P(M).
\]
前者称为平凡拓扑（indiscrete topology），后者称为离散拓扑（discrete topology）。

\begin{remark}[拓扑不是距离]
拓扑不要求我们给每一对点指定一个数值距离。它只记录哪些集合可以作为“开放的局部区域”。在某些空间中可以由距离生成拓扑，但拓扑本身比某个具体距离更一般。
\end{remark}

\subsection{标准欧氏拓扑}

对$x\in\R^d$和$r>0$，定义以$x$为中心、半径为$r$的开球：
\[
    B_r(x)
    :=
    \{y\in\R^d:\|y-x\|<r\}.
\]
若采用欧氏范数，则
\[
    \|y-x\|^2
    =
    \sum_{i=1}^d(y^i-x^i)^2.
\]
因此也可以写成
\[
    B_r(x)
    =
    \left\{
    y\in\R^d:
    \sum_{i=1}^d(y^i-x^i)^2<r^2
    \right\}.
\]

\begin{definition}[标准欧氏拓扑]
$\R^d$上的标准拓扑定义为
\[
    \mathcal O_{\mathrm{std}}
    =
    \left\{
    U\subseteq\R^d:
    \forall x\in U,\ \exists r>0,\ B_r(x)\subseteq U
    \right\}.
\]
\end{definition}

直观地说：
\[
    U\text{ 是开集}
    \iff
    \text{$U$中的每个点都能放进一个完全留在$U$内的小球。}
\]

\begin{remark}[球不是唯一的局部形状]
在$\R^2$中，小正方形、小矩形和小圆盘都可以用来生成同一个标准拓扑。定义时使用开球最简洁；“正方形邻域”的图像只是对开球思想的可视化，不改变拓扑定义。
\end{remark}

\subsection{邻域与局部性质}

\begin{definition}[邻域]
设$(M,\mathcal O)$是拓扑空间，$p\in M$。若$N\subseteq M$包含某个含$p$的开集，则称$N$是$p$的一个邻域。
\end{definition}

注意，邻域本身不一定是开集。只要存在开集$U$满足
\[
    p\in U\subseteq N,
\]
则$N$就是$p$的邻域。

\begin{definition}[局部性质]
若某个性质在点$p$的某个邻域中成立，则称该性质在$p$处局部成立。
\end{definition}

流形的局部欧氏性、连续性的局部判别以及局部坐标化，都会反复使用这个观点：
\[
    \text{先在每个点附近处理，再研究这些局部处理能否拼成整体结构。}
\]

%——————————————————————————————————%

\section{子空间拓扑与积拓扑}

\subsection{子空间拓扑}

\begin{definition}[子空间拓扑]
设$(M,\mathcal O)$是拓扑空间，$N\subseteq M$。定义
\[
    \mathcal O_N
    :=
    \{U\cap N:U\in\mathcal O\}.
\]
称$\mathcal O_N$为$N$从$M$继承的子空间拓扑，$(N,\mathcal O_N)$称为$M$的子空间。
\end{definition}

一个集合$A\subseteq N$在子空间$N$中开，当且仅当存在$M$中的开集$U$满足
\[
    A=U\cap N.
\]
这里的$A$不必是$M$中的开集；它只需要是环境空间开集切到$N$上所得的相对开集。

\begin{proposition}[子空间拓扑确实是拓扑]
$\mathcal O_N$满足拓扑的三条公理。
\end{proposition}

验证的关键是集合运算：
\[
    \bigcup_{\alpha\in A}(U_\alpha\cap N)
    =
    \left(\bigcup_{\alpha\in A}U_\alpha\right)\cap N,
\]
以及有限交
\[
    \bigcap_{i=1}^k(U_i\cap N)
    =
    \left(\bigcap_{i=1}^kU_i\right)\cap N.
\]
右侧中的并集和有限交仍属于$\mathcal O$，所以左侧属于$\mathcal O_N$。

\subsection{相对开集与相对闭集}

令
\[
    N=[-1,1]\subset\R.
\]
集合
\[
    (0,1]
\]
不是$\R$中的开集，但它是$N$中的开集，因为
\[
    (0,1]
    =
    (0,2)\cap[-1,1].
\]
因此必须区分：
\[
    \boxed{
    (0,1]\text{ 不在$\R$中开，但在$[-1,1]$中开。}
    }
\]

同理，$C\subseteq N$在$N$中闭，当且仅当存在$M$中的闭集$F$使
\[
    C=F\cap N.
\]
“开”或“闭”总是相对于当前讨论的拓扑空间而言，不能脱离环境空间直接判断。

\begin{remark}[曲面与环境空间]
球面
\[
    S^2=\{x\in\R^3:\|x\|=1\}
\]
不是$\R^3$中的开集，但它由子空间拓扑自身成为一个拓扑空间。以后说$S^2$中的开集，指的是从$\R^3$的开集切到$S^2$上得到的相对开集。
\end{remark}

\subsection{圆周的两种拓扑来源}

把圆周作为$\R^2$的子集：
\[
    S^1
    =
    \{(x,y)\in\R^2:x^2+y^2=1\},
\]
则自然使用子空间拓扑：
\[
    \mathcal O_{S^1}
    =
    \{U\cap S^1:U\text{ 是$\R^2$中的开集}\}.
\]

圆周也可以由实数轴每隔$2\pi$识别一次得到。定义等价关系
\[
    \theta\sim\varphi
    \iff
    \theta-\varphi\in2\pi\Z.
\]
商集合为
\[
    \R/(2\pi\Z).
\]
自然投影为
\[
    q:\R\to\R/(2\pi\Z),
    \qquad
    q(\theta)=[\theta].
\]

\begin{definition}[商拓扑]
若$q:X\to X/{\sim}$是自然投影，则商空间上的商拓扑定义为：
\[
    A\subseteq X/{\sim}\text{ 开}
    \iff
    q^{-1}(A)\subseteq X\text{ 开}.
\]
\end{definition}

映射
\[
    \bar q:\R/(2\pi\Z)\to S^1,
    \qquad
    \bar q([\theta])=(\cos\theta,\sin\theta)
\]
是一个同胚。这给出：
\[
    \R/(2\pi\Z)
    \cong_{\mathrm{top}}
    S^1.
\]

这里必须注意：
\[
    \R/(2\pi\Z)
\]
作为集合并不会自动带有拓扑；商拓扑是使上述“绕一圈识别”的构造成为连续空间的必要结构。

\subsection{积拓扑}

设$(A,\mathcal O_A)$和$(B,\mathcal O_B)$是拓扑空间。集合层面的积空间是$A\times B$。

\begin{definition}[积拓扑]
$A\times B$上的积拓扑是由所有矩形集合
\[
    U\times V,
    \qquad
    U\in\mathcal O_A,\quad V\in\mathcal O_B,
\]
生成的最小拓扑。
\end{definition}

等价地，$W\subseteq A\times B$开，当且仅当对任意$(a,b)\in W$，存在$U\in\mathcal O_A$和$V\in\mathcal O_B$，使
\[
    a\in U,\qquad b\in V,\qquad U\times V\subseteq W.
\]
直观上，每个点周围都能找到一个完全落在$W$内的小矩形。

\begin{proposition}[欧氏空间的有限积]
标准欧氏空间满足
\[
    \R^d
    \cong
    \underbrace{\R\times\cdots\times\R}_{d\text{ 次}},
\]
并且$\R^d$的标准欧氏拓扑与有限积拓扑一致。
\end{proposition}

有限积中的坐标投影
\[
    \operatorname{pr}_1:A\times B\to A,
    \qquad
    \operatorname{pr}_2:A\times B\to B
\]
是连续映射。以后纤维丛的投影、局部平凡化以及状态空间
\[
    \SE(3)^N\times\R^{3K}
\]
都会使用积空间结构。

%——————————————————————————————————%

\section{度量、闭集与局部结构}

\subsection{度量空间}

\begin{definition}[度量]
集合$M$上的函数
\[
    d:M\times M\to[0,\infty)
\]
若满足：
\begin{enumerate}
    \item 非负性：
    \[
        d(p,q)\ge0;
    \]
    \item 正定性：
    \[
        d(p,q)=0\iff p=q;
    \]
    \item 对称性：
    \[
        d(p,q)=d(q,p);
    \]
    \item 三角不等式：
    \[
        d(p,r)\le d(p,q)+d(q,r);
    \]
\end{enumerate}
则称$d$为$M$上的一个度量，$(M,d)$称为度量空间。
\end{definition}

度量生成开球：
\[
    B_r(p)=\{q\in M:d(p,q)<r\}.
\]
由所有开球生成的拓扑称为度量拓扑。

\begin{theorem}[欧氏拓扑由距离诱导]
在$\R^d$中取欧氏距离
\[
    d(x,y)=\|x-y\|,
\]
由开球生成的度量拓扑恰好是标准欧氏拓扑。
\end{theorem}

拓扑语言比距离语言更一般：一个空间可能有多个度量生成同一个拓扑，也可能根本没有我们当前需要的自然度量。研究流形的第一步通常只需要拓扑结构，度量和内积将在后面的微分几何中重新出现。

\subsection{闭集}

\begin{definition}[闭集]
设$(M,\mathcal O)$是拓扑空间。集合$C\subseteq M$称为闭集，当且仅当补集
\[
    M\setminus C
\]
是开集。
\end{definition}

空集和全空间总是既开又闭：
\[
    \varnothing,\ M
\]
都是clopen集合。既开又闭并不意味着集合“既有内部又没有边界”，它只是相对于当前拓扑同时满足开集和闭集的定义。

\begin{lemma}[开闭对偶]
任意并的补集满足
\[
    M\setminus\bigcup_{\alpha\in A}U_\alpha
    =
    \bigcap_{\alpha\in A}(M\setminus U_\alpha),
\]
任意交的补集满足
\[
    M\setminus\bigcap_{\alpha\in A}U_\alpha
    =
    \bigcup_{\alpha\in A}(M\setminus U_\alpha).
\]
\end{lemma}

这解释了拓扑公理的对偶形式：
\[
\begin{array}{c|c}
\text{开集}&\text{闭集}\\
\hline
\text{任意并}&\text{任意交}\\
\text{有限交}&\text{有限并}
\end{array}
\]

\subsection{内部、闭包与边界}

\begin{definition}[内部]
对$A\subseteq M$，$A$的内部记为$\operatorname{Int}A$，定义为包含于$A$的最大开集：
\[
    \operatorname{Int}A
    =
    \bigcup\{U\subseteq A:U\text{ 是开集}\}.
\]
\end{definition}

\begin{definition}[闭包]
对$A\subseteq M$，$A$的闭包记为$\overline A$，定义为包含$A$的最小闭集：
\[
    \overline A
    =
    \bigcap\{C\supseteq A:C\text{ 是闭集}\}.
\]
\end{definition}

\begin{definition}[边界]
集合$A\subseteq M$的边界定义为
\[
    \partial A
    =
    \overline A\setminus\operatorname{Int}A.
\]
\end{definition}

对任意$A\subseteq M$，
\[
    \operatorname{Int}A\subseteq A\subseteq\overline A.
\]
在$\R$的标准拓扑中：
\[
    \operatorname{Int}[0,1]=(0,1),
    \qquad
    \overline{(0,1)}=[0,1],
    \qquad
    \partial[0,1]=\{0,1\}.
\]

边界、内部和闭包都依赖于环境拓扑。对$A\subseteq N\subseteq M$，在$N$中取内部和边界，可能与在$M$中取出的结果不同。

%——————————————————————————————————%

\section{序列与收敛}

\subsection{序列}

\begin{definition}[序列]
拓扑空间$M$中的序列是一个映射
\[
    q:\N\to M.
\]
通常记为
\[
    \{q(n)\}_{n\in\N}
    \quad\text{或}\quad
    q_1,q_2,\ldots.
\]
\end{definition}

序列可以表示迭代优化中的估计：
\[
    x_0,x_1,x_2,\ldots,
\]
也可以表示传感器采样、轨迹离散化或数值算法产生的状态序列。

\subsection{拓扑收敛}

\begin{definition}[序列收敛]
序列$q(n)$在拓扑空间$M$中收敛到$p\in M$，若对$p$的每个开邻域$U$，都存在$N\in\N$，使得
\[
    n\ge N
    \Longrightarrow
    q(n)\in U.
\]
记为
\[
    q(n)\to p.
\]
\end{definition}

用量词完整写出：
\[
    \forall U\in\mathcal O,\ p\in U,
    \quad
    \exists N\in\N,\ \forall n\ge N,\ q(n)\in U.
\]

在$\R^d$中，这就是熟悉的$\varepsilon$语言：
\[
    q(n)\to p
    \iff
    \forall\varepsilon>0,\ \exists N,\ 
    n\ge N\Longrightarrow\|q(n)-p\|<\varepsilon.
\]

\subsection{普通欧氏空间中的收敛}

例如
\[
    q(n)=1+\frac1n
\]
在$\R$的标准拓扑中收敛到$1$：
\[
    q(n)\to1.
\]
它不是最终恒定的，但从某一项开始会进入$1$的任意小开区间。

\begin{remark}[最终性质]
收敛只要求序列最终进入每一个给定邻域，不要求序列最终恒定。数值优化中，估计量通常逐步接近极限，而不是在某一轮之后完全不再变化。
\end{remark}

\subsection{离散拓扑与平凡拓扑}

这两个例子必须严格区分。

\begin{definition}[离散拓扑]
若
\[
    \mathcal O=\P(M),
\]
则称$M$带有离散拓扑。
\end{definition}

离散拓扑中每个单点集都是开集。于是：
\[
    q(n)\to p
    \iff
    \text{$q(n)$最终恒等于$p$}.
\]
证明只需取开邻域$\{p\}$。收敛要求存在$N$使
\[
    n\ge N\Longrightarrow q(n)\in\{p\},
\]
也就是$q(n)=p$。

\begin{definition}[平凡拓扑]
若
\[
    \mathcal O=\{\varnothing,M\},
\]
则称$M$带有平凡拓扑。
\end{definition}

平凡拓扑中每个点的唯一非空开邻域是$M$本身，因此任意序列都收敛到每一个点：
\[
    \forall q:\N\to M,\ \forall p\in M,\qquad q(n)\to p.
\]

\begin{center}
\begin{tabular}{c|c}
\text{拓扑}&\text{序列收敛行为}\\
\hline
$\P(M)$\text{（离散拓扑）}
&
\text{收敛当且仅当最终恒定}\\
$\{\varnothing,M\}$\text{（平凡拓扑）}
&
\text{每个序列收敛到每个点}
\end{tabular}
\end{center}

这说明拓扑会实质性地改变“极限”的含义。

%——————————————————————————————————%

\section{连续映射与同胚}

\subsection{连续映射}

\begin{definition}[拓扑连续]
设$(M,\mathcal O_M)$和$(N,\mathcal O_N)$是拓扑空间。映射
\[
    f:M\to N
\]
称为连续映射，当且仅当对$N$中的每个开集$V$，都有
\[
    f^{-1}(V)
    =
    \{p\in M:f(p)\in V\}
\]
是$M$中的开集。
\end{definition}

连续性的定义关注陪域开集的原像：
\[
    \boxed{
    V\subseteq N\text{ 开}
    \Longrightarrow
    f^{-1}(V)\subseteq M\text{ 开}.
    }
\]

一般而言，$f$连续并不保证定义域开集$U$的像$f(U)$在陪域中开。原像与像的方向不能混淆。

\begin{theorem}[连续复合]
若$f:M\to N$和$g:N\to P$连续，则复合映射
\[
    g\circ f:M\to P
\]
连续。
\end{theorem}

\begin{proof}
设$W\subseteq P$开。因为$g$连续，$g^{-1}(W)$在$N$中开；因为$f$连续，
\[
    f^{-1}(g^{-1}(W))
    =
    (g\circ f)^{-1}(W)
\]
在$M$中开。
\end{proof}

\begin{proposition}[连续性的局部性]
若$\{U_\alpha\}$是$M$的开覆盖，且每个限制映射
\[
    f\vert_{U_\alpha}:U_\alpha\to N
\]
连续，则$f:M\to N$连续。
\end{proposition}

这说明连续性可以先在局部区域验证，再利用开覆盖拼接为整体性质。流形上的局部坐标、地图拼接和局部观测模型都会使用这一原则。

\subsection{极端拓扑下的连续映射}

若陪域$N$带平凡拓扑：
\[
    \mathcal O_N=\{\varnothing,N\},
\]
则任意映射$f:M\to N$连续，因为
\[
    f^{-1}(\varnothing)=\varnothing,
    \qquad
    f^{-1}(N)=M.
\]

若定义域$M$带离散拓扑：
\[
    \mathcal O_M=\P(M),
\]
则任意映射$f:M\to N$连续，因为每个原像都是$M$的子集，因而都是开集。

\subsection{序列判别}

\begin{theorem}[度量空间中的序列判别]
设$M,N$是度量空间。映射$f:M\to N$连续，当且仅当对任意序列$q(n)\to p$，都有
\[
    f(q(n))\to f(p).
\]
\end{theorem}

在一般拓扑空间中，序列未必足以检测所有开集；上述“连续当且仅当保持序列极限”的结论需要度量空间或更一般的第一可数条件。对本笔记中的欧氏空间、常见流形和SLAM状态空间局部模型，序列判别通常足够直观。

\begin{corollary}[观测链连续]
若状态到中间变量、投影、相机内参和残差映射均连续，则它们的复合观测模型连续。
\end{corollary}

\subsection{同胚}

\begin{definition}[同胚]
若映射$f:M\to N$是双射，且$f$与$f^{-1}$都连续，则称$f$为同胚。记为
\[
    f:M\xrightarrow{\cong_{\mathrm{top}}}N
\]
或
\[
    M\cong_{\mathrm{top}}N.
\]
\end{definition}

同胚空间在拓扑意义下不可区分：它们具有相同的开集结构、邻近关系、连通性和紧致性。它们可能在几何形状、坐标表示或嵌入环境中不同，但拓扑结构相同。

\begin{proposition}[坐标图思想]
流形上的一张坐标图本质上是一个局部邻域$U$与$\R^d$中开集之间的同胚：
\[
    U\cong_{\mathrm{top}}x(U)\subseteq\R^d.
\]
\end{proposition}

这正是第四章“局部欧氏性”的拓扑核心。

%——————————————————————————————————%

\section{分离公理：从\texorpdfstring{$T_1$}{T1}到\texorpdfstring{$T_3$}{T3}}

“分离”可以理解为：拓扑中的开集能够把不同对象区分到什么程度。条件越强，需要同时满足的开集关系越严格；这里的“强”是逻辑意义上的强，而不是给拓扑空间赋予某种数值大小。

\subsection{\(T_1\)空间}

\begin{definition}[$T_1$空间]
拓扑空间$M$称为$T_1$空间，若对任意不同点$p\neq q$，存在开集$U$满足
\[
    q\in U,\qquad p\notin U.
\]
\end{definition}

交换$p,q$后也必须成立，因此$T_1$空间中的两个不同点可以分别被开集排除。

\begin{theorem}[$T_1$与单点闭集]
拓扑空间$M$是$T_1$空间，当且仅当对任意$p\in M$，单点集
\[
    \{p\}
\]
是闭集。
\end{theorem}

\begin{example}[有限补集拓扑：$T_1$但不是Hausdorff]
令$X$是无限集合，并定义
\[
    \mathcal O_{\mathrm{cof}}
    =\{\varnothing\}\cup
    \{U\subseteq X:X\setminus U\text{是有限集}\}.
\]
这称为$X$上的有限补集拓扑。对任意$p\neq q$，集合
\[
    U=X\setminus\{p\}
\]
是开集，且$q\in U$、$p\notin U$；交换$p,q$同样成立，所以它是$T_1$空间。

但任意两个非空开集$U,V$都满足$U\cap V\neq\varnothing$。否则
\[
    X=(X\setminus U)\cup(X\setminus V)
\]
会成为有限集，与$X$无限矛盾。因此，两个不同点虽然可以分别排除对方，却不能拥有不相交的开邻域。这个例子说明
\[
    T_1\not\Rightarrow T_2.
\]
\end{example}

\subsection{Hausdorff空间}

\begin{definition}[Hausdorff空间]
拓扑空间$M$称为Hausdorff空间，也称为$T_2$空间，若对任意不同点$p\neq q$，存在开集$U,V$满足
\[
    p\in U,\qquad q\in V,\qquad U\cap V=\varnothing.
\]
\end{definition}

例如在$\R$的标准拓扑中，取$p=0$和$q=2$，令
\[
    U=\left(-\frac12,\frac12\right),
    \qquad
    V=\left(\frac32,\frac52\right).
\]
则$p\in U$、$q\in V$且$U\cap V=\varnothing$。与有限补集拓扑相比，Hausdorff条件要求两个点的邻域必须同时选出，并且完全不相交。

\subsection{分离公理的层级}

为了看清这些条件的强弱关系，先补充最弱的一层。

\begin{definition}[$T_0$空间]
拓扑空间$M$称为$T_0$空间，若对任意不同点$p\neq q$，至少存在一个开集能够包含其中一点而不包含另一点。
\end{definition}

再进一步，称拓扑空间$M$为正则空间，若对任意点$p$和不含$p$的闭集$F$，存在开集$U,V$满足
\[
    p\in U,
    \qquad
    F\subseteq V,
    \qquad
    U\cap V=\varnothing.
\]
本文采用常见约定：同时满足正则性和$T_1$性的空间称为$T_3$空间。

\begin{figure}[htbp]
\centering
\begin{tikzpicture}[
    >=stealth,
    level/.style={
        draw=structurecolor!70,
        fill=structurecolor!8,
        rounded corners=2pt,
        align=center,
        text width=2.55cm,
        minimum height=0.95cm,
        inner sep=4pt
    },
    implication/.style={->, line width=0.7pt, draw=structurecolor!85}
]
    \node[level] (t3) at (-4.5,0)
        {$T_3$\\[-2pt]\small 点与闭集};
    \node[level] (t2) at (-1.5,0)
        {$T_2$\\[-2pt]\small 两点不交邻域};
    \node[level] (t1) at (1.5,0)
        {$T_1$\\[-2pt]\small 逐点排除};
    \node[level] (t0) at (4.5,0)
        {$T_0$\\[-2pt]\small 至少一方可区别};
    \draw[implication] (t3.east) -- (t2.west);
    \draw[implication] (t2.east) -- (t1.west);
    \draw[implication] (t1.east) -- (t0.west);
\end{tikzpicture}
\caption{分离公理的强弱关系}
\label{fig:separation-axioms}
\end{figure}

在$T_3$的定义中取$F=\{q\}$即可得到$T_2$：因为$T_1$空间的单点集是闭集，所以点与另一个点可以被不交开集分开。因此
\[
    T_3\Longrightarrow T_2\Longrightarrow T_1\Longrightarrow T_0.
\]
这些蕴含一般不能反向成立；有限补集拓扑已经给出了$T_1$而非$T_2$的例子，而两点集$\{0,1\}$上的拓扑
\[
    \{\varnothing,\{1\},\{0,1\}\}
\]
是$T_0$但不是$T_1$。

\begin{example}[度量空间是$T_3$]
设$(M,d)$是度量空间，$p\notin F$且$F$是闭集。由于$M\setminus F$是开集，存在$r>0$使得
\[
    B_d(p,r)\cap F=\varnothing.
\]
取
\[
    U=B_d\left(p,\frac r3\right),
    \qquad
    V=M\setminus\overline{B_d\left(p,\frac r3\right)}.
\]
则$U,V$都是开集，$p\in U$，$F\subseteq V$且$U\cap V=\varnothing$。度量空间又是$T_1$空间，所以它是$T_3$空间；特别地，$\R$的标准拓扑同时满足$T_3$、$T_2$和$T_1$。

这个例子揭示了三层要求的差别：$T_1$只关心“能否排除另一个点”，$T_2$要求“两个点能否拥有互不相交的邻域”，而$T_3$进一步要求“一个点能否和整个闭集分开”。
\end{example}

\begin{theorem}[Hausdorff空间中的极限唯一性]
若$M$是Hausdorff空间，且
\[
    q(n)\to p,\qquad q(n)\to r,
\]
则
\[
    p=r.
\]
\end{theorem}

\begin{proof}
反设$p\neq r$。取不相交开邻域$U\ni p$和$V\ni r$。由收敛性，从某个指标开始$q(n)\in U$，从另一个指标开始$q(n)\in V$。取二者中较大的指标后，得到某个$q(n)$同时属于$U$和$V$，这与$U\cap V=\varnothing$矛盾。
\end{proof}

\begin{remark}[流形为什么通常要求Hausdorff]
如果两个不同的物理位置没有不相交邻域，就很难把它们作为可区分的状态。Hausdorff性保证极限唯一，也排除了许多不适合作为物理状态空间的病态拓扑。
\end{remark}

%——————————————————————————————————%

\section{紧致性}

\subsection{开覆盖与有限子覆盖}

\begin{definition}[开覆盖]
设$\mathcal U=\{U_\alpha\}_{\alpha\in A}$是一族开集。若
\[
    M=\bigcup_{\alpha\in A}U_\alpha,
\]
则称$\mathcal U$为$M$的一个开覆盖。
\end{definition}

\begin{definition}[紧致空间]
拓扑空间$M$称为紧致空间，若$M$的任意开覆盖都有有限子覆盖。
\end{definition}

也就是说，若
\[
    M=\bigcup_{\alpha\in A}U_\alpha,
    \qquad U_\alpha\text{ 开},
\]
则存在有限多个指标$\alpha_1,\ldots,\alpha_k$，使得
\[
    M
    =
    U_{\alpha_1}\cup\cdots\cup U_{\alpha_k}.
\]

关键词是：
\[
    \boxed{
    \text{任意给定的开覆盖中，都能挑出有限个成员仍覆盖全体。}
    }
\]
紧致性不是“存在一个有限开覆盖”，而是“每一个开覆盖都可以被有限控制”。

\subsection{基本紧致性定理}

\begin{theorem}[Heine--Borel定理]
对$K\subseteq\R^d$，
\[
    K\text{ 紧致}
    \iff
    K\text{ 闭且有界}.
\]
\end{theorem}

\begin{remark}[定理的适用范围]
Heine--Borel定理是欧氏空间中的特殊结论。一般拓扑空间或一般度量空间中，“闭且有界”未必推出紧致；紧致性的定义始终是开覆盖定义。
\end{remark}

\begin{theorem}[闭子集紧致]
若$K$是紧致空间，且$C\subseteq K$在$K$中闭，则$C$紧致。
\end{theorem}

\begin{theorem}[连续像紧致]
若$K$紧致，$f:K\to N$连续，则$f(K)$紧致。
\end{theorem}

\begin{corollary}[Hausdorff空间中的紧子集闭]
若$N$是Hausdorff空间，$K\subseteq N$紧致，则$K$在$N$中闭。
\end{corollary}

这些定理把“有限控制”传递到子空间和连续变换之后。它们也是后面证明极值存在和理解紧流形的重要工具。

\subsection{极值存在}

\begin{theorem}[Weierstrass定理]
若$K\subseteq\R^d$紧致，且
\[
    f:K\to\R
\]
连续，则$f$在$K$上取得最大值和最小值。即存在$p_{\min},p_{\max}\in K$，使
\[
    f(p_{\min})
    =
    \min_{p\in K}f(p),
    \qquad
    f(p_{\max})
    =
    \max_{p\in K}f(p).
\]
\end{theorem}

\begin{corollary}[有界搜索区域]
若SLAM目标函数连续，且搜索区域是非空的闭有界集合，则全局最小值存在。
\end{corollary}

\begin{remark}[存在性不等于可计算性]
紧致性保证极值存在，但不保证数值算法能找到全局极小点，也不保证局部线性化、初值或噪声条件良好。它解决的是“最优解是否存在”，不是“算法如何找到它”。
\end{remark}

\subsection{紧致性与状态估计}

设$\mathcal X$是状态空间，$\mathcal F\subseteq\mathcal X$是由先验边界、传感器范围或物理约束得到的可行集合。若$\mathcal F$紧致，则连续代价函数在$\mathcal F$上有最小值。

在欧氏状态空间中，常用方式是限制到闭有界区域：
\[
    \mathcal F
    =
    \{x:g_i(x)\ge0,\ i=1,\ldots,m\}
    \cap
    \{x:\|x\|\le R\}.
\]
在流形状态空间中，则需要先确认相应集合的拓扑和闭性；不能直接把“参数向量有界”当作整个流形状态有界。

%——————————————————————————————————%

\section{可数性、仿紧性与分割统一}

\subsection{拓扑基与第二可数性}

\begin{definition}[拓扑基]
设$\mathcal B\subseteq\mathcal O$。若对任意开集$U\in\mathcal O$和任意$p\in U$，存在$B\in\mathcal B$满足
\[
    p\in B\subseteq U,
\]
则称$\mathcal B$为$\mathcal O$的一个拓扑基。
\end{definition}

\begin{definition}[第二可数]
若拓扑空间存在一个可数拓扑基，则称其为第二可数空间。
\end{definition}

第二可数性意味着整个拓扑可以由可数多个基本开集生成。它不是说空间只有可数多个点，而是说描述开集结构所需的基本模板可以取为可数族。

\begin{remark}[欧氏空间的第二可数性]
$\R^d$的标准拓扑有可数基，例如以有理数点为中心、以正有理数为半径的开球族：
\[
    \mathcal B
    =
    \{B_r(x):x\in\Q^d,\ r\in\Q_{>0}\}.
\]
\end{remark}

\subsection{覆盖、细化与局部有限}

\begin{definition}[覆盖的细化]
设$\mathcal U$和$\mathcal V$是$M$的开覆盖。若对每个$V\in\mathcal V$，存在$U\in\mathcal U$满足
\[
    V\subseteq U,
\]
则称$\mathcal V$是$\mathcal U$的一个开细化。
\end{definition}

\begin{definition}[局部有限族]
集合族$\mathcal V$称为局部有限，若对每个$p\in M$，存在$p$的开邻域$W_p$，使得$W_p$至多与$\mathcal V$中的有限多个集合相交：
\[
    \#\{V\in\mathcal V:W_p\cap V\neq\varnothing\}<\infty.
\]
\end{definition}

\begin{definition}[仿紧空间]
拓扑空间$M$称为仿紧空间（paracompact space），若它的每个开覆盖都有一个局部有限的开细化。
\end{definition}

局部有限的意义是：虽然整体上可能有无穷多个局部区域，但在每一个具体点附近，真正发生作用的区域只有有限多个。这样局部构造的求和或拼接才不会在一个点处累积无穷多项。

\begin{theorem}[紧空间仿紧]
每个紧致空间都是仿紧空间。
\end{theorem}

\begin{theorem}[流形的仿紧性]
Hausdorff、第二可数且局部欧氏的拓扑流形是仿紧空间。
\end{theorem}

这个结论是第四章中采用Hausdorff和第二可数条件的重要原因之一：它保证局部坐标、局部函数和局部几何对象可以通过良好控制的覆盖拼成全局对象。

\subsection{分割统一}

\begin{definition}[支集]
连续函数$\varphi:M\to\R$的支集定义为
\[
    \operatorname{supp}\varphi
    =
    \overline{\{p\in M:\varphi(p)\neq0\}}.
\]
\end{definition}

\begin{definition}[分割统一]
给定开覆盖
\[
    \mathcal U=\{U_i\}_{i\in I},
\]
一个从属于$\mathcal U$的分割统一，是一族连续函数
\[
    \varphi_i:M\to[0,1],
\]
满足：
\begin{enumerate}
    \item $\{\operatorname{supp}\varphi_i\}_{i\in I}$局部有限；
    \item
    \[
        \operatorname{supp}\varphi_i\subseteq U_i;
    \]
    \item 对所有$p\in M$，
    \[
        \sum_{i\in I}\varphi_i(p)=1.
    \]
\end{enumerate}
\end{definition}

必须要求支集满足
\[
    \operatorname{supp}\varphi_i\subseteq U_i,
\]
而不能只要求
\[
    \varphi_i(p)\neq0\Longrightarrow p\in U_i.
\]
后一个条件只控制非零集合，不控制其闭包；标准定义需要整个支集都从属于$U_i$。

\begin{theorem}[分割统一的存在性]
在适当的正则性条件下，特别是在光滑流形上，任意开覆盖都有从属于它的光滑分割统一。
\end{theorem}

若$f_i$是在$U_i$上定义的局部函数，则先把$\varphi_i f_i$在$U_i$外按零延拓。由于
\[
    \operatorname{supp}\varphi_i\subseteq U_i,
\]
这个延拓在相应条件下保持连续。于是可以用
\[
    f(p)=\sum_{i\in I}\varphi_i(p)f_i(p)
\]
构造全局函数。局部有限性保证每个$p$附近只有有限项真正参与求和。

\begin{remark}[局部到整体]
分割统一是“局部定义、全局拼接”的标准工具。它并不把不同局部坐标简单地当成同一套坐标，而是用权重函数在重叠区域中平滑组合。
\end{remark}

%——————————————————————————————————%

\section{连通性与道路连通}

\subsection{连通空间}

\begin{definition}[连通]
拓扑空间$M$称为连通，若不存在两个非空、不交的开集$A,B$满足
\[
    M=A\cup B,
    \qquad
    A\cap B=\varnothing.
\]
\end{definition}

这样的表达称为把$M$分裂成两个开集。连通性表示空间不能被拓扑地拆成两个彼此分离的非空部分。

\begin{theorem}[连通与clopen集合]
拓扑空间$M$连通，当且仅当$M$中唯一既开又闭的集合是
\[
    \varnothing,\qquad M.
\]
\end{theorem}

\begin{theorem}[连续像保持连通]
若$M$连通且$f:M\to N$连续，则$f(M)$连通。
\end{theorem}

\subsection{道路连通}

\begin{definition}[道路连通]
若对任意$p,q\in M$，都存在连续映射
\[
    \gamma:[0,1]\to M
\]
满足
\[
    \gamma(0)=p,\qquad \gamma(1)=q,
\]
则称$M$道路连通。
\end{definition}

在$\R^d$中，连接$p,q$的直线段为
\[
    \gamma(t)=(1-t)p+tq,
    \qquad t\in[0,1].
\]
因此$\R^d$道路连通。

\begin{theorem}[道路连通推出连通]
任一道路连通空间都是连通空间。
\end{theorem}

\begin{remark}[反向一般不成立]
连通性只要求空间不能被分裂成两个开闭部分；道路连通性要求任意两点之间存在一条连续路径。因此道路连通是更强的条件，不能把二者视为同义词。
\end{remark}

\subsection{道路分支}

定义点$p\in M$的道路分支为所有能由连续路径与$p$连接的点组成的集合。道路分支将空间按道路连通关系分成若干部分。

\begin{remark}[连通分支与道路分支]
在一般拓扑空间中，连通分支和道路分支不一定相同；在许多常见的流形和局部道路连通空间中，二者的行为更规整。对SLAM而言，若状态空间存在多个道路分支，连续运动模型无法从一个分支直接运动到另一个分支。
\end{remark}

\subsection{位姿空间的道路连通性}

\begin{theorem}[$\SO(3)$道路连通]
三维旋转空间$\SO(3)$是道路连通的。
\end{theorem}

\begin{theorem}[$\SE(3)$道路连通]
三维位姿空间$\SE(3)$是道路连通的。
\end{theorem}

\begin{corollary}[连续位姿路径]
任意两个三维刚体位姿之间存在连续运动路径。
\end{corollary}

这里的道路连通性只说明路径存在，不说明路径唯一，也不说明沿路径的速度、能量或优化代价具有特殊性质。后续微分几何和李群理论会进一步研究这些局部和代数结构。

%——————————————————————————————————%

\section{同伦与基本群}

\subsection{路径同伦}

\begin{definition}[同伦]
设$f,g:X\to Y$是连续映射。若存在连续映射
\[
    H:X\times[0,1]\to Y
\]
满足
\[
    H(x,0)=f(x),
    \qquad
    H(x,1)=g(x),
    \qquad
    \forall x\in X,
\]
则称$f$与$g$同伦，记为
\[
    f\simeq g.
\]
\end{definition}

对两条路径
\[
    \gamma,\delta:[0,1]\to M,
\]
同伦映射为
\[
    H:[0,1]\times[0,1]\to M.
\]
参数$t$表示路径上的位置，参数$s$表示形变过程：
\[
    H(t,0)=\gamma(t),
    \qquad
    H(t,1)=\delta(t).
\]

若路径端点固定，则要求对所有$s\in[0,1]$：
\[
    H(0,s)=\gamma(0)=\delta(0),
    \qquad
    H(1,s)=\gamma(1)=\delta(1).
\]
这种同伦称为保持端点的路径同伦。

\subsection{环路与基点}

\begin{definition}[基于$p$的环路]
给定$p\in M$，基于$p$的环路是连续映射
\[
    \gamma:[0,1]\to M
\]
满足
\[
    \gamma(0)=\gamma(1)=p.
\]
\end{definition}

若$\gamma,\delta$是基于$p$的环路，基点保持同伦是连续映射
\[
    H:[0,1]\times[0,1]\to M
\]
满足
\[
    H(t,0)=\gamma(t),
    \qquad
    H(t,1)=\delta(t),
\]
并且
\[
    H(0,s)=p,
    \qquad
    H(1,s)=p,
    \qquad
    \forall s\in[0,1].
\]

最后两条条件说明形变过程中，环路的起点和终点始终钉在基点$p$。如果不固定基点，得到的是另一种路径同伦关系，不能直接用于定义基于$p$的基本群。

\subsection{路径拼接}

若路径$\gamma,\delta:[0,1]\to M$满足
\[
    \gamma(1)=\delta(0),
\]
则定义拼接路径
\[
    (\gamma*\delta)(t)
    =
    \begin{cases}
        \gamma(2t),&0\le t\le\frac12,\\[4pt]
        \delta(2t-1),&\frac12\le t\le1.
    \end{cases}
\]
前半段走完$\gamma$，后半段走完$\delta$。在拼接点处连续性来自
\[
    \gamma(1)=\delta(0).
\]

路径反向定义为
\[
    \gamma^{-1}(t)=\gamma(1-t).
\]
常值路径记为
\[
    c_p(t)=p.
\]

\subsection{基本群}

\begin{definition}[基本群]
固定基点$p\in M$。所有基于$p$的环路在基点保持同伦下的等价类构成集合
\[
    \pi_1(M,p).
\]
其群运算由路径拼接诱导：
\[
    [\gamma]\cdot[\delta]
    =
    [\gamma*\delta].
\]
\end{definition}

在适当的同伦等价意义下：
\[
    [c_p]
\]
是单位元，
\[
    [\gamma^{-1}]
\]
是$[\gamma]$的逆元。路径拼接在参数化层面未必严格结合，但在同伦等价类层面结合，因此基本群运算是良定义的。

\begin{theorem}[基本群的典型例子]
\[
    \pi_1(\R^d,p)=\{[c_p]\},
    \qquad
    \pi_1(S^2,p)=\{[c_p]\},
    \qquad
    \pi_1(S^1,p)\cong\Z.
\]
\end{theorem}

$\R^d$和$S^2$中的每个环路都可以连续缩成常值环路；圆周中的环路则有绕行次数这一整数不变量：
\[
    \cdots,-2,-1,0,1,2,\cdots.
\]
绕圆一圈、两圈或反向绕行，不能在圆周内部连续缩成一点。

\begin{remark}[基本群与全局结构]
局部上，$S^1$和$\R$都像一条直线；但全局上$S^1$存在无法收缩的闭合路径。这说明局部欧氏性不能完全决定空间的整体拓扑，图册之外还需要研究连通、环路和同伦等全局性质。
\end{remark}

\subsection{与SLAM的联系}

SLAM状态空间中的连续轨迹、闭环运动和周期变量都可能携带全局拓扑信息。局部坐标可以描述当前位置附近的增量，但不能自动消除：

\begin{itemize}
    \item 状态空间的多个道路分支；
    \item 周期变量绕行一周后的全局识别；
    \item 闭合路径是否可以连续缩回起点；
    \item 局部参数化在全局上是否存在奇点。
\end{itemize}

因此，在进入拓扑流形和李群之前，需要保留一个基本判断：
\[
    \text{局部坐标适合做局部计算，但全局拓扑决定哪些运动和表示真正可能。}
\]

%——————————————————————————————————%

\section{本章小结：从拓扑到流形}

本章的核心内容可以压缩为以下几条：

\begin{enumerate}
    \item 拓扑是在集合$M$上指定一族开集$\mathcal O$，满足空集、全空间、任意并和有限交公理。
    \item 拓扑公理的不对称性是：
    \[
        \text{任意并仍开，但只有有限交仍开。}
    \]
    \item 子空间$N\subseteq M$使用相对拓扑：
    \[
        \mathcal O_N=\{U\cap N:U\in\mathcal O_M\}.
    \]
    \item 积空间使用矩形基$U\times V$；有限积给出：
    \[
        \R^d\cong\R\times\cdots\times\R.
    \]
    \item 序列$q(n)$收敛到$p$，表示它最终进入$p$的每一个开邻域。离散拓扑中收敛等价于最终恒定；平凡拓扑中每个序列收敛到每个点。
    \item 连续映射由开集原像定义；同胚是双向连续的双射。
    \item Hausdorff性保证不同点有不相交邻域，并保证序列极限唯一。
    \item 紧致性表示任意开覆盖都有有限子覆盖；在$\R^d$中等价于闭且有界，并由此得到连续函数的极值存在性。
    \item 仿紧性提供局部有限细化；分割统一要求
    \[
        \operatorname{supp}\varphi_i\subseteq U_i,
        \qquad
        \sum_i\varphi_i=1.
    \]
    \item 道路连通推出连通，但连通不必推出道路连通；基本群记录基点保持同伦下的环路结构。
\end{enumerate}

进入第四章后，拓扑空间$M$还要满足局部欧氏性、Hausdorff性和第二可数性：
\[
    \forall p\in M,\quad
    \exists U\subseteq M,\quad
    p\in U,\quad
    U\cong_{\mathrm{top}}x(U)\subseteq\R^d,
\]
其中$U$是$M$中的开集，$x(U)$是$\R^d$中的开集。这就是拓扑流形的入口。

后续层次关系为
\[
    \text{拓扑空间}
    \longrightarrow
    \text{拓扑流形}
    \longrightarrow
    \text{可微流形}
    \longrightarrow
    \text{李群}
    \longrightarrow
    \text{流形上的SLAM优化}.
\]
